% !TEX root = ../../main_docker.tex
% Grundbegriffe

% ---------------------------------------------------------------------
%  Docker -- Software-Virtualisierung mit Containern
% ---------------------------------------------------------------------
\section{Grundbegriffe}
\label{sec:docker}

\SW{Wozu} Plattform zur \T{Software-Virtualisierung}: Software wird in
standardisierte, wiederverwendbare Einheiten verpackt, die \emph{alles}
enthalten, was zum Ausführen nötig ist. Diese Einheiten laufen auf nahezu
jedem Rechner direkt \ra portabel, skalierbar, unkompliziert.
Das Problem, das sie löst: „Bei mir läuft's!“ -- Abhängigkeiten,
Versionskonflikte, Installationsprozeduren.

\SW{Abgrenzung} Docker virtualisiert \emph{nur den Linux-Kernel}, nicht
die ganze Maschine. Deshalb gilt immer \T{Container-OS $=$ Host-OS} --
anders als bei einer virtuellen Maschine. Auf Windows/macOS braucht Docker
deshalb eine VM (WSL). Die Programme laufen direkt auf dem Host-Kernel;
in \R{ps ax} auf dem Host sieht man den Container-Prozess.

\subsection{Image vs.\ Container}
\begin{ztab}{@{}lXX@{}}
 & Image & Container \\ \midrule
Was & Bauplan, alles Nötige zum Ausführen & ein \emph{ausgeführtes} Image \\
Dateisystem & \T{read-only}, unveränderlich & \T{read-write} Overlay über dem Image \\
Anzahl & eines & beliebig viele aus einem Image \\
Platz & volle Grösse & nur die \emph{Änderungen} gegenüber dem Image \\
Lebensdauer & bis \R{rmi} & bis \R{rm}, überlebt das Stoppen \\
\end{ztab}

\begin{formel}[Der Satz, der alles trägt]
Das Ausführen eines \emph{Images} erzeugt einen \emph{Container} --
nie umgekehrt. Änderungen zur Laufzeit (Logfiles, DB-Einträge,
nachinstallierte Pakete) landen im Overlay des Containers, das Image
bleibt unberührt.
\end{formel}

\subsection{Host und Kapselung}
\SW{Host} der Rechner mit dem OS, auf dem die Container laufen (hier
Raspberry Pi OS). Er stellt nur die grundlegendsten OS-Funktionen bereit:
nacktes Dateisystem, Prozesse, Sockets. \emph{Nicht} Software, Daten,
User/Gruppen oder IP-Ports -- das muss der Container selbst mitbringen.

\SW{Kapselung} Aus dem Container heraus normalerweise kein Zugriff auf das
Dateisystem des Hosts oder anderer Container; gilt auch für User und Ports.
Der Host dagegen kann Container und Images ändern und löschen.
Die zwei Löcher in dieser Wand sind \R{docker cp} und \T{Volumes}.
